% Pursuing Stacks (cotes 134-2 to 134-8): archive cote 134-7, site page 100 (PDF page 101), typescript page 542.
% typescript. Draft compared: ps6.tex:5038-5106.
% AI-transcribed by Claude Opus 5.5 (claude-opus-5-5, adaptive thinking, effort high), Message Batches API, batch ps4 (msgbatch_012Aoc1NuskATPsyBeKSYK7J), 2026-11-12; not yet checked by eye.
\scanpage{134-7/100}
\mspage{542}

14.10.

129.\quad I have not been clear enough in yesterday's notes, when intro-%
ducing the linearization functor (11)
\[ L:\ \ U(k)\longrightarrow Uab(k) \]
for a suitable notion of ``bundles'' and ``linear bundles'', that it is by
no means automatic that such a functor (even when it's existence is gran%
ted) will induce a linearization functor
\[ (16)\qquad LH.\ :\ Hot_{o}(k)\longrightarrow Hotab_{o}(k) \]
for the corresponding derived categories. In other words, it is by no
means clear (and we hav'nt even tried yet to prove in the $U_{W}$ - set-up
of ``unipotent bundles'') that if
\[ (17)\qquad X_{*}\longrightarrow Y_{*} \]
is a weak equivalence of semisimplicial objects of U(k), that the
corresponding map
\[ (17')\qquad L(X_{*})\longrightarrow L(Y_{*}) \]
of semisimplicial objects in Uab(k) will again be a weak equivalence.
At any rate, for this statement even to make sense, we'll have to make
clear what notion of weak equivalence we are taking for \inserted{maps between} semisimplicial
objects of Uab(k), so as to get a localized category Hotab$_{\unsure{\bullet|0}}$(k). The
first obvious choice of course that comes to\xout{x} mind is to take the notion
of weak equivalence for the corresponding semisimplicial sheaves of sets\hand{,}
which presumably is going to be the right one. Thus, the existence of
a ``\inserted{total}homology'' or ``linearization'' functor (16) may be viewed as the
``schematic'' analog of W.H.C. Whitehead's theorem in the discrete set-up,
which we had been puzzling about already in section 92 (in connection
with replacing the test category \hand{$\Delta$} by a more or less arbitrary small
category). In the case of semisimplicial W-bundles or V-bundles, un-%
der some extra assumption presumably such as $X_{o}$=e or $X_{\unsure{\bullet|0}}$=$X_{1}$=e on the
semisimplicial objects we are working with, I do expect that the line-%
rization functor transforms \xout{quasi-isomorphisms} weak equivalences into
weak equivalences\hand{,} and that the converse holds true too. This seems to
me to be the third main test (besides the two considered on previous
page 541) about a notion of ``bundle'' being suited for develop\deleted{p}ing a
theory of schematic homotopy types.

This suggests that when working with more general semisimplicial
sheaves over k, such as semisimplicial schemes say, it may be useful

\margin{100}
